\section{Relation to the literature}
\label{sec:literature}

The contribution here is deliberately source-specific.  This paper corrects and globally recharacterizes the model of \cite{ShyStenbacka2005}; the contribution is not a new general theorem that competition reduces outsourcing, that investment games generate asymmetric equilibria, or that no-loss pricing is novel.

Subsequent outsourcing research studies several distinct mechanisms.  These include real options \cite{AlvarezStenbacka2007}, organizational forms \cite{ShyStenbacka2012} and \cite{StenbackaTombak2012}, and procurement or strategic sourcing \cite{Rosar2017}.  Other work studies component structure \cite{ChenChen2020}, quality choice \cite{Miyamoto2021}, and strategic interactions between make-or-buy choices and supplier investment \cite{Dai2026}.  Dai's mechanism is distinct: the organizational choice changes an upstream firm's investment incentive and thereby affects rivals' input costs.  We do not propose a new make-or-buy mechanism of that kind.  Most importantly for the competition comparative static, \cite{Konig2010} studies two outsourcing motives.  In his reversed specification, where outsourcing saves marginal rather than fixed costs, intensified competition reduces outsourcing, and he explicitly contrasts that result with \cite{ShyStenbacka2005}.  Our result is narrower still: equation (14) in the published Shy--Stenbacka model itself has the opposite sign from the one reported there.

The asymmetric-equilibrium result also belongs to a well-developed strategic-investment literature.  Cost-reducing investment followed by Cournot competition can produce polarized or asymmetric outcomes and, in some specifications, endogenous exit \cite{Amir2000,AmirHalmenschlagerJin2011}.  The symmetry-breaking class in \cite{AmirGarciaKnauff2010} emphasizes strategic substitutes together with payoff nonconcavity, while \cite{BuehlerSchmutzler2008} studies asymmetric investment incentives generated by endogenous vertical integration.  Other work establishes multiple asymmetric investment equilibria \cite{EckertKlumppSu2017} and piecewise innovation responses \cite{LamantiaPezzinoTramontana2018}.  These results absorb any generic novelty claim about asymmetry or multiplicity.  What is specific here is the corrected correspondence inside the published sourcing model: own reduced payoff remains globally strictly concave, and multiplicity can arise even while both downstream firms remain active.  Separately, when the sourcing cap is large enough, a downstream exit threshold creates the positive-slope best-response piece that invalidates a global strategic-substitutes reading of Proposition 5.

The Hotelling result is likewise not a claim of a new general spatial-pricing phenomenon.  Heterogeneous-cost spatial competition is standard \cite{Vogel2008}.  In a different posted-price/discount model, \cite{GillThanassoulis2016} explicitly restrict both list and discount prices to be no lower than marginal cost; that restriction provides context for, but does not derive, the auxiliary price domain used here.  Our literal-source result is instead that sufficiently asymmetric sourcing histories move the published Hotelling continuation outside its interior branch and generate multiple pure corner continuations.  The \(p_j\ge c_j\) analysis is a separate robustness exercise used to show that the published symmetric sourcing candidate can also fail under a transparent no-loss strategy restriction.

These distinctions matter for interpretation.  The results establish a correction of the published model and a complete pure Stage-I re-characterization of its Cournot source-duopoly case.  They should not be read as establishing generic comparative statics for arbitrary demand, monitoring technologies, or mixed-strategy equilibrium sets.
